\documentclass[../../../include/open-logic-section]{subfiles}

\begin{document}

\olfileid{sth}{choice}{countablechoice}
\olsection{गणनीय निवड}

निवड किंवा सुक्रमण अजिबात न वापरता पुढील निष्कर्ष सहज सिद्ध करता येतो:

\begin{lem}[$\Zminus$ मध्ये]
प्रत्येक सांत संचाला निवडफलन असते.
\end{lem}

\begin{proof}
$a = \{b_1, \ldots, b_n\}$ असू द्या; सदस्यांची पुनरावृत्ती करू नका. सुलभतेसाठी प्रत्येक $b_i
\neq \emptyset$ आहे असे समजा. मग $c_1 \in b_1, \ldots, c_n \in b_n$ असलेल्या वस्तू $c_1, \ldots, c_n$ आहेत. आता \olref[z][pairs]{prop:pairsconsequences} नुसार $\{\langle b_1,
c_1\rangle , \ldots, \langle b_n, c_n\rangle\}$ हा संच अस्तित्वात आहे; आणि हे $a$ साठीचे निवडफलन आहे. रिक्त सदस्य वगळता येतो; रिक्त कुटुंबासाठी रिक्त फलन पुरेसे आहे.
\end{proof}

पण अनंत संचांचा विचार केल्याबरोबर गोष्टी अधिक अस्पष्ट होतात. उदाहरणार्थ, वरील विधानाचा हा “किमान” विस्तार पाहा:

\begin{defish}
\emph{गणनीय निवड.} प्रत्येक \emph{गणनीय} संचाला निवडफलन असते.
\end{defish}

हे निवडीचे एक विशेषप्रकरण आहे. हे तत्त्व वापरत असल्याची स्पष्ट जाणीव नसतानाही ते बऱ्याचदा वापरले गेले होते, असे दिसते. पुढील दोन सुंदर उदाहरणे पाहा.\footnote{ही उदाहरणे \citet[\S9.4]{Potter2004} आणि लूका इन्कुर्वाती यांच्याकडून घेतली आहेत.}

\begin{ex}
एक स्वाभाविक विचार असा: कोणत्याही संच $A$ साठी $\cardle{\omega}{A}$ किंवा एखाद्या $n \in \omega$ साठी $\cardeq{A}{n}$. प्रत्येक संच सांत किंवा अनंत असतो, ही सहज कल्पना मांडण्याचा हा एक मार्ग आहे. कँटर आणि इतर अनेक गणितज्ञांनी पुराव्याशिवाय हा दावा केला. आपण सावध असल्यामुळे \olref[cardinals][classing]{generalinfinitycharacter} मध्ये तो सिद्ध केला. पण त्या पुराव्यात कोणत्याही संच $A$ चे सुक्रमण करता येते आणि म्हणून $\card{A}$ अस्तित्वात असल्याची खात्री आहे, असे मानून आपण $\ZFC$ मध्ये काम केले. म्हणजे आपण निवड स्पष्टपणे गृहीत धरली होती.

खरे तर \citet{Dedekind1888} यांनी या दाव्याचा स्वतःचा पुरावा पुढीलप्रमाणे दिला होता:

\begin{thm}[$\Zminus + \text{गणनीय निवड}$ मध्ये]
कोणत्याही $A$ साठी $\cardle{\omega}{A}$ किंवा एखाद्या $n \in \omega$ साठी $\cardeq{A}{n}$.
\end{thm}

\begin{proof}
प्रत्येक $n \in \omega$ साठी $\cardneq{A}{n}$ असे समजा. मग विशेषतः प्रत्येक $n < \omega$ साठी नेमके $\cardexpo{2}{n}$ घटक असलेला उपसंच $A_n \subseteq A$ आहे. $A_0, A_1, A_2,
\ldots$ हा अनुक्रम वापरून प्रत्येक $n$ साठी पुढील व्याख्या करू:
\[
	B_n = A_n \setminus \bigcup_{i < n} A_i.
\]
शून्यव्या निर्देशांकासाठी आधीचे युतीकरण रिक्त आहे आणि पहिला उपसंच अरिक्त आहे. धन निर्देशांकांसाठी पुढील मोजणी पाहा:
\begin{align*}
	\card{\bigcup_{i < n}A_i}
	&\leq \card{A_0} + \card{A_1} + \ldots + \card{A_{n-1}}\\
	&=1 + 2 + \ldots + 2^{n-1}\\
	& = 2^n - 1\\
	& < 2^n = \card{A_n}
\end{align*}
म्हणून प्रत्येक $B_n$ मध्ये किमान एक घटक $c_n$ आहे. शिवाय $B_n$ हे संच जोडीजोडीने वियुक्त आहेत; त्यामुळे $c_n = c_m$ असल्यास $n = m$. पण प्रत्येक $c_n
\in A$. म्हणून $f(n) = c_n$ हे फलन $\omega \to A$ असे एकास-एक फलन आहे.
\end{proof}
\noindent
गणनीय निवड वापरली असल्याची डेडेकिंडने नोंद केली नाही. पण \emph{तुम्हाला} तिचा वापर दिसला का? पुन्हा पाहा. (खरेच: \emph{पुन्हा पाहा}.)

या पुराव्यात गणनीय निवड दोनदा वापरली. $A_0$, $A_1$, $A_2$, \dots\@ या संचांचा अनुक्रम मिळवण्यासाठी ती एकदा वापरली. त्यानंतर प्रत्येक $B_n$ मधून $c_n$ हे घटक निवडण्यासाठी ती पुन्हा वापरली. शिवाय येथे कोणत्याही प्रकारची निवड न वापरता हा निष्कर्ष सिद्ध करता येत नाही. \citet[p.~138]{Cohen1966} यांनी निवडीचे कोणतेही रूप नसल्यास हा निष्कर्ष अपयशी ठरतो, असे सिद्ध केले. म्हणजे $\omega$ शी \emph{अतुलनीय} असलेले संच आहेत, ही बाब $\ZF$ शी सुसंगत आहे.
\end{ex}

\begin{ex}
\citeyear{Cantor1878} मध्ये कँटरने म्हटले होते की गणनीय संचांचे गणनीय युतीकरण गणनीय असते. त्याने पुरावा दिला नाही; कदाचित तो पुरावा स्वयंस्पष्ट आहे, असे त्याला वाटले असावे. आपण सावध असल्यामुळे \olref[card-arithmetic][simp]{kappaunionkappasize} मध्ये या निष्कर्षाचे अधिक सामान्य रूप सिद्ध केले. पण आपल्या पुराव्यात निवड स्पष्टपणे गृहीत धरली होती. आणि कमी सामान्य निष्कर्षाचा खालील पुरावाही गणनीय निवड वापरतो.

\begin{thm}[$\Zminus + \text{गणनीय निवड}$ मध्ये]
प्रत्येक $n \in \omega$ साठी $A_n$ गणनीय असेल, तर $\bigcup_{n <
\omega} A_n$ गणनीय असतो.
\end{thm}

\begin{proof}
सर्व सदस्य रिक्त असतील, तर त्यांचे युतीकरण रिक्त आणि गणनीय आहे. अन्यथा रिक्त सदस्य वगळून आणि आवश्यकतेनुसार अरिक्त सदस्यांची अनुक्रमात पुनरावृत्ती करून, सामान्यतेला बाधा न आणता प्रत्येक $A_n \neq \emptyset$ आहे असे मानता येते. मग प्रत्येक $n \in \omega$ साठी !!a{surjection} $f_n \colon \omega
\to A_n$ आहे. $f(m, n) = f_n(m)$ घेऊन $f \colon \omega \times \omega \to \bigcup_{n <
\omega} A_n$ परिभाषित करा. $\omega
\times \omega$ गणनीय आहे (\olref[sfr][siz][zigzag]{natsquaredenumerable}) आणि $f$ हे !!a{surjection} आहे, म्हणून निष्कर्ष मिळतो.
\end{proof}
\noindent
गणनीय निवडीचा वापर दिसला का? $f_0$, $f_1$, $f_2$, \dots\footnote{\olref[card-arithmetic][simp]{kappaunionkappasize} मध्ये “प्रत्येक $\beta \in \cardfont{a}$ साठी !!a{injection} $f_\beta$ निश्चित करा” अशी सूचना देताना निवडीचा असाच वापर झाला होता.} हा फलनांचा अनुक्रम निवडण्यासाठी ती वापरली आहे. पुन्हा, कोणतेही निवडतत्त्व नसल्यास हा निष्कर्ष अपयशी ठरू शकतो. विशेषतः \citet{FefermanLevy1963} यांनी असे सिद्ध केले की गणनीय संचांच्या गणनीय युतीकरणाचा संचांक $\beth_1$ असतो, ही बाब $\ZF$ शी सुसंगत आहे. पण याच मुद्द्याचे रसेल यांनी केलेले अधिक मजेशीर वर्णन पाहा:
\begin{quote}
  एक कोट्यधीश जेव्हा बुटांची जोडी विकत घेत असे, तेव्हाच मोज्यांचीही जोडी घेत असे, इतर वेळी कधीच नाही. दोन्ही गोष्टी विकत घेण्याची त्याला इतकी हौस होती की शेवटी त्याच्याकडे बुटांच्या $\aleph_0$ जोड्या आणि मोज्यांच्या $\aleph_0$ जोड्या झाल्या\dots\@ बुटांत उजवा आणि डावा असा फरक करता येतो. त्यामुळे प्रत्येक जोडीतून एक बूट निवडता येतो: सर्व उजवे किंवा सर्व डावे बूट निवडू शकतो. पण मोज्यांच्या बाबतीत असे कोणतेही निवडतत्त्व सुचत नाही; आणि गुणक स्वयंसिद्धक [म्हणजे प्रत्यक्षात निवड] गृहीत धरल्याशिवाय, प्रत्येक जोडीतून एक मोजा असलेला वर्ग आहे याची खात्री देता येत नाही.
  \citep[p.~126]{Russell1919}
\end{quote}
म्हणजे मोज्यांच्या गणनीय संख्येने जोड्या असतील, तर मोजेही फक्त गणनीय संख्येने असतात, हे सिद्ध करण्यासाठी काही प्रकारची निवड लागते. गणनीय निवड अशा युतीकरणाची गणनीयता सिद्ध करण्यास पुरेशी आहे; पण निवडतत्त्वांचा अभाव असताना गणनीय संचांचे गणनीय युतीकरण गणनीय असण्यात अपयश येऊ शकते. येथे गणनीय निवडीशी उलटी सममूल्यता सिद्ध केलेली नाही.
\end{ex}

यातून धडा असा मिळतो की वापरणाऱ्यांना फारशी जाणीव नसतानाही गणनीय निवड वारंवार वापरली गेली. तात्त्विक प्रश्न असा आहे: गणनीय निवडीचे \emph{समर्थन} कसे करावे?

सहज समर्थनाचा एखादा प्रयत्न सांत वेळेत अनंत क्रियांचा क्रम पूर्ण करण्याच्या कल्पनेचा आधार घेऊ शकतो. पहिली निवड $\nicefrac{1}{2}$ मिनिटात, दुसरी $\nicefrac{1}{4}$ मिनिटात, \dots, आणि $n$ वी निवड $\nicefrac{1}{2^n}$ मिनिटात, \dots\@ केली असे समजा. मग $1$ मिनिटात निवडींचा $\omega$-अनुक्रम पूर्ण करून निवडफलन परिभाषित केले असेल.

पण असा विचारप्रयोग खरोखर काय सांगू शकतो? प्रथम, “निवड करणे” ही कल्पना अगदी शब्दशः घेतल्यावर तो अवलंबून आहे. दुसरे म्हणजे त्यात गणित सत्तामीमांसक शक्यतेशी बांधले जात असल्यासारखे दिसते.

अधिक महत्त्वाचे असे की यामुळे \emph{केवळ} गणनीय निवडीऐवजी निवडीच्या \emph{पूर्ण रूपाचे} समर्थन मिळणार नाही. कारण \emph{प्रत्येक} संचाला निवडफलन हवे असेल, तर “कोणत्याही क्रमसूचक लांबीचा अनंत क्रियाक्रम” पूर्ण करता आला पाहिजे. स्पष्ट सांगायचे तर ही कल्पना हास्यास्पद आहे.

\end{document}
